\begin{solapartat}
La reacció nuclear del ${}^{210}_{84}\mathrm{Po}$, serà:
\[ {}^{210}_{84}\mathrm{Po} \rightarrow {}^{4}_{2}\mathrm{He} + {}^{206}_{82}\mathrm{Pb}\ \text{\punts{0,5}} \]
També considerem vàlida la resposta on en lloc del He i posem $\alpha$

Podem escriure la llei de desintegració com:
\[ N(t) = N_0\,\mathrm{e}^{-\lambda\,t} \]
Per altre banda:
\begin{align*}
N(t = t_{1/2}) &= \frac{N_0}{2} = N_0\,\mathrm{e}^{-\lambda\,t_{1/2}}\ \text{\punts{0,2}} \Rightarrow \mathrm{e}^{\lambda\,t_{1/2}} = 2 \Rightarrow\\
\lambda &= \frac{\ln 2}{t_{1/2}} = 5,02{\cdot}10^{-3}\,\mathrm{dies^{-1}} = 5,81{\cdot}10^{-8}\,\mathrm{s^{-1}}\ \text{\punts{0,3}}
\end{align*}
\end{solapartat}

\begin{solapartat}
La llei de desintegració també la podem escriure:
\[ \frac{m(t)}{m_A}N_A = \frac{m_0}{m_A}N_A\,\mathrm{e}^{-\lambda\,t}\ \text{\punts{0,1}} \Rightarrow m(t) = m_0\,\mathrm{e}^{-\lambda\,t}\ \text{\punts{0,2}} \]
Per tant:
\[ m(t = 20\ \text{dies}) = 5\ \mathrm{mg}\ \mathrm{e}^{-\,5,02\cdot 10^{-3}\cdot\mathrm{dies}^{-1}\cdot 20\ \mathrm{dies}} = 4,52\ \mathrm{mg}\ \text{\punts{0,7}} \]
O sigui ens quedaran: 4,52~mg
\end{solapartat}
